\documentclass[11pt]{article}

\usepackage[margin=1in]{geometry}
\usepackage[T1]{fontenc}
\usepackage{lmodern}

\usepackage{amsmath, amssymb}

\usepackage{setspace}

\usepackage[most]{tcolorbox}
\usepackage{xcolor}

\tcbset{
  enhanced,
  boxrule=0pt,
  frame hidden,
  breakable
}

\usepackage{fancyhdr}
\pagestyle{fancy}
\fancyhf{}
\setlength{\headheight}{52pt}
\renewcommand{\headrulewidth}{0pt}

\newcommand{\Course}{Math 4610 Spring 2026}
\newcommand{\Instructor}{Discussion 6}
\newcommand{\AssignmentType}{Discussion} % or Discussion
\newcommand{\AssignmentNumber}{6}

\newcommand{\Contributors}{
Marks, Stephen\\
}

\fancyhead[L]{\Course\\ \Instructor}
\fancyhead[R]{\Contributors}

\newcommand{\ProblemDivider}[1]{%
  \vspace{1.0em}
  \noindent\textbf{#1}\par
  \vspace{0.35em}
}

\definecolor{problemBack}{RGB}{235,242,250}
\definecolor{problemFrame}{RGB}{70,110,160}
\definecolor{exerciseGray}{gray}{0.96}
\definecolor{exerciseGrayFrame}{gray}{0.45}

\newtcolorbox{problemBox}{
  colback=exerciseGray,
  colframe=exerciseGrayFrame,
  arc=2mm,
  left=8pt,right=8pt,top=8pt,bottom=8pt
}

\newtcolorbox{solutionBox}{
  colback=white,
  colframe=white,
  arc=0mm,
  left=0pt,right=0pt,top=0pt,bottom=0pt
}

\newcommand{\ProblemAnnounce}{\noindent}
\newcommand{\SolutionAnnounce}{\noindent\textbf{Solution.}\par\medskip}

\newcommand{\Exercise}[1]{\textbf{Exercise~#1.}\;}

\newcommand{\R}{\mathbb{R}}
\newcommand{\C}{\mathbb{C}}

\begin{document}

\begin{center}
  {\Large \textbf{\AssignmentType~\AssignmentNumber}}
\end{center}
\vspace{0.75em}

\ProblemDivider{1}

\begin{problemBox}
\ProblemAnnounce
\Exercise{1}
Recall that
\[
R_{\theta}(x,y) = (x\cos\theta - y\sin\theta,\; x\sin\theta + y\cos\theta)
\]
is the \(\theta\)-rotation of the plane \(R_{\theta} \in \mathcal{L}(\C^2)\).

\begin{enumerate}
    \item[(a)] Fix \(\theta \in (0,2\pi)\). Show that the eigenvalues of \(R_{\theta}\) are \(e^{i\theta} = \cos\theta + i\sin\theta\) and \(e^{-i\theta} = \cos\theta - i\sin\theta\).
    \item[(b)] Find the corresponding eigenvectors.
\end{enumerate}
\end{problemBox}

\begin{solutionBox}
\SolutionAnnounce
\doublespacing

\begin{enumerate}
    \item[(a)] We want \(\lambda \in \C\) and nonzero \((x,y) \in \C^2\) such that
    \[
    R_{\theta}(x,y) = \lambda (x,y).
    \]
    This gives the system:
    \[
    \begin{cases}
        x\cos\theta - y\sin\theta = \lambda x \\[4pt]
        x\sin\theta + y\cos\theta = \lambda y
    \end{cases}
    \]
    Then:
    \[
    \begin{cases}
        (\cos\theta - \lambda)x - (\sin\theta) y = 0 \\[4pt]
        (\sin\theta)x + (\cos\theta - \lambda) y = 0
    \end{cases}
    \]
    If \(\sin\theta \neq 0\), we have
    \[
    y = \frac{\cos\theta - \lambda}{\sin\theta}\, x.
    \]
    By substituting into the second equation:
    \[
    \sin\theta \cdot x + (\cos\theta - \lambda) \cdot \frac{\cos\theta - \lambda}{\sin\theta}\, x = 0.
    \]
    Multiply by \(\sin\theta\):
    \[
    \sin^2\theta\, x + (\cos\theta - \lambda)^2 x = 0.
    \]
    Since \(x \neq 0\),
    \[
    \sin^2\theta + (\cos\theta - \lambda)^2 = 0.
    \]
    By expanding \((\cos\theta - \lambda)^2 = \cos^2\theta - 2\lambda\cos\theta + \lambda^2\) and using \(\sin^2\theta + \cos^2\theta = 1\):
    \[
    1 - 2\lambda\cos\theta + \lambda^2 = 0.
    \]
    Thus
    \[
    \lambda^2 - 2\lambda\cos\theta + 1 = 0.
    \]
    Solving the quadratic:
    \[
    \lambda = \cos\theta \pm \sqrt{\cos^2\theta - 1}
          = \cos\theta \pm i\sin\theta.
    \]
    Then
    \[
    \lambda_1 = e^{i\theta} = \cos\theta + i\sin\theta,\quad
    \lambda_2 = e^{-i\theta} = \cos\theta - i\sin\theta.
    \]

    \item[(b)] For \(\lambda_1 = e^{i\theta} = \cos\theta + i\sin\theta\):
    \[
    (\cos\theta - (\cos\theta + i\sin\theta))x - (\sin\theta) y = 0
    \]
    \[
    (-i\sin\theta)x - \sin\theta\, y = 0.
    \]
    Assuming \(\sin\theta \neq 0\), divide by \(-\sin\theta\):
    \[
    i x + y = 0 \quad\Longrightarrow\quad y = -i x.
    \]
    Choose \(x = 1\), then \(y = -i\). So an eigenvector is \((1, -i)\).

    For \(\lambda_2 = e^{-i\theta} = \cos\theta - i\sin\theta\):
    \[
    (\cos\theta - (\cos\theta - i\sin\theta))x - (\sin\theta) y = 0
    \]
    \[
    (i\sin\theta)x - \sin\theta\, y = 0.
    \]
    Divide by \(\sin\theta\):
    \[
    i x - y = 0 \quad\Longrightarrow\quad y = i x.
    \]
    Choose \(x = 1\), then \(y = i\). So an eigenvector is \((1, i)\).

    Thus the eigenvectors are \((1,-i)\) for \(e^{i\theta}\) and \((1,i)\) for \(e^{-i\theta}\).
\end{enumerate}

\end{solutionBox}
\ProblemDivider{2}

\begin{problemBox}
\ProblemAnnounce
\Exercise{2}
If \(\mathcal{L}(\R^2,\R)\) stands for the operators from \(\R^2\) to \(\R\), show that \(\mathcal{L}(\R^2,\R)\) and \(\R^2\) are isomorphic vector spaces. This is a continuation of Problem 6 in Homework 2.

[Hint: Define \(G:\mathcal{L}(\R^2,\R) \longrightarrow \R^2\) as follows: If \(T:\R^2 \longrightarrow \R\) is linear, set
\[
G(T) := (T(1,0), T(0,1)).
\]
Notice that \(G(T)\) is the vector in \(\R^2\) whose entries are the images under \(T\) of the elements in the canonical basis \(\{(1,0),(0,1)\}\). Then show that \(G\) is linear, 1-1 and onto. For example, if \(T(x,y) = 2x + 3y\), then \(G(T) = (T(1,0),T(0,1)) = (2,3)\).]
\end{problemBox}

\begin{solutionBox}
\SolutionAnnounce
\doublespacing

Define \(G: \mathcal{L}(\R^2,\R) \to \R^2\) by
\[
G(T) = (T(1,0), T(0,1)).
\]

\noindent \textbf{Linearity of \(G\):}
Let \(S,T \in \mathcal{L}(\R^2,\R)\) and \(c \in \R\). Then
\[
G(S+T) = ((S+T)(1,0), (S+T)(0,1)) = (S(1,0)+T(1,0), S(0,1)+T(0,1))
\]
\[
= (S(1,0), S(0,1)) + (T(1,0), T(0,1)) = G(S) + G(T).
\]
Also,
\[
G(cT) = ((cT)(1,0), (cT)(0,1)) = (c T(1,0), c T(0,1))
= c (T(1,0), T(0,1)) = c G(T).
\]
Thus \(G\) is linear.

\noindent \textbf{Injectivity:}
Suppose \(G(T) = (0,0)\). Then \(T(1,0) = 0\) and \(T(0,1) = 0\). For any \((x,y) \in \R^2\),
\[
T(x,y) = T(x(1,0) + y(0,1)) = x T(1,0) + y T(0,1) = x\cdot 0 + y\cdot 0 = 0.
\]
Thus \(T\) is the zero operator, so \(\ker G = \{0\}\) and \(G\) is injective.

\noindent \textbf{Surjectivity:}
Let \((a,b) \in \R^2\). Define \(T: \R^2 \to \R\) by \(T(x,y) = ax + by\). Then \(T\) is linear,
\[
T(1,0) = a,\quad T(0,1) = b,
\]
so \(G(T) = (a,b)\). Thus \(G\) is surjective.

Since \(G\) is linear, injective, and surjective, it is an isomorphism. Therefore
\[
\mathcal{L}(\R^2,\R) \cong \R^2.
\]

\end{solutionBox}

\end{document}
